\section{Уроки 8--9: камера, кадры и калибровка}

\subsection{Цели уроков}
\begin{itemize}
\item Получать видеопоток с дрона в формате OpenCV.
\item Сохранять кадры и подготавливать данные для CV-задач.
\item Выполнять калибровку камеры и сохранять коэффициенты в \texttt{data.yml}.
\end{itemize}

\subsection{Опорные листинги}
\lstinputlisting[style=GeoskanStyle,language=pythoncustom,caption={Запуск и остановка видеопотока}]{../pioneer-python-example/camera_examples/camera_stream.py}
\lstinputlisting[style=GeoskanStyle,language=pythoncustom,caption={Получение кадров с камеры}]{../pioneer-python-example/camera_examples/get_frames_from_camera.py}
\lstinputlisting[style=GeoskanStyle,language=pythoncustom,caption={Калибровка камеры по шахматной доске}]{../pioneer-python-example/camera_examples/camera_calibration.py}

\subsection{Методика калибровки}
\begin{enumerate}
\item Снять 10--15 кадров шахматной доски под разными углами.
\item Убедиться, что углы доски распознаются на большинстве кадров.
\item Сохранить матрицу камеры и коэффициенты дисторсии.
\item Повторно использовать \texttt{data.yml} в задачах ArUco.
\end{enumerate}

\begin{taskbox}[Минимальный набор упражнений]
\begin{enumerate}
\item Подготовьте набор из 12 кадров для калибровки.
\item Сохраните коэффициенты и проверьте, что файл \texttt{data.yml} создан.
\item Объясните, почему без калибровки оценка расстояния до маркера нестабильна.
\end{enumerate}
\end{taskbox}

\subsection{Критерии успешности}
\begin{enumerate}
\item Скрипт калибровки завершается без ошибок OpenCV.
\item В \texttt{data.yml} записаны \texttt{mtx} и \texttt{dist}.
\item При повторном запуске CV-скриптов файл корректно загружается.
\end{enumerate}
